% from ega4-11-fr.tex
\oldpage[IV-3]{152}
fixé, l'intersection des idéaux
$\widehat{\mathfrak m}_\alpha'^{\,n}\widehat{A'}_{\widehat{\mathfrak m}'_\alpha}$
pour $n>0$ est réduite à $0$ dans l'anneau local noethérien
$\widehat{A'}_{\widehat{\mathfrak m}'_\alpha}$. D'autre part les
$\widehat{\mathfrak m}'_\alpha$ sont les idéaux maximaux de $\widehat{A'}$, donc
l'homomorphisme canonique
$\widehat{A'}\to\prod_\alpha\widehat{A'}_{\widehat{\mathfrak m}'_\alpha}$ est
injectif (Bourbaki, \emph{Alg. comm.}, chap.~II, \S~3, n$^\circ$~3, cor.~2 du
th.~1), et comme par hypothèse $\widehat u:\widehat A\to\widehat{A'}$ est aussi
injectif, cela achève la démonstration dans le cas général. La dernière
assertion résulte de ce que $\widehat A$ est un $A$-module fidèlement plat
($\mathbf 0_{\mathrm I},7.3.5$) et $\widehat{A'}=A'\otimes_A\widehat A$ puisque
$A'$ est par hypothèse un $A$-module de type fini (Bourbaki, \emph{Alg. comm.},
chap.~III, \S~3, n$^\circ$~4, th.~3 et chap.~IV, \S~2, n$^\circ$~5, cor.~3 de la
prop.~9).
\end{proof}

\begin{proposition}[11.5.3]
\label{IV.11.5.3-fr}
Soient $f:X\to Y$ un morphisme localement de présentation finie,
$\mathcal F$ un $\mathcal O_X$-Module quasi-cohérent de présentation finie, $x$
un point de $X$, $y=f(x)$, $g=(\psi,\theta):Y'\to Y$ un morphisme propre de
présentation finie. Supposons que~:
\begin{enumerate}[label=(\roman*)]
  \item l'homomorphisme $\theta_y:\mathcal O_y\to
  (g_*\mathcal O_{Y'})_y$ soit injectif~;
  \item pour tout $x'\in X'=X\times_Y Y'$ dont la projection dans $X$ est $x$,
  $\mathcal F'=\mathcal F\otimes_Y Y'$ soit $Y'$-plat au point $x'$.
\end{enumerate}
Alors $\mathcal F$ est $f$-plat au point $x$.
\end{proposition}

\begin{proof}
La question étant locale sur $X$, on peut supposer $f$ de présentation finie.
Soit $p:X'\to X$ la projection canonique. Comme
$f'=f_{(Y')}:X'\to Y'$ est de présentation finie (1.6.2) et que $\mathcal F'$
est un $\mathcal O_{X'}$-Module de présentation finie
$(\mathbf 0_{\mathrm I},5.2.5)$, il résulte de (11.3.1) que l'ensemble $U'$ des
points de $X'$ où $\mathcal F'$ est $f'$-plat est ouvert. Comme $U'$ contient
$p^{-1}(x)$ par hypothèse, et que $p$ est propre, donc fermé, $U'$ contient un
ensemble de la forme $p^{-1}(U)$, où $U$ est un voisinage de $x$. Remplaçant $X$
par $U$, on peut donc supposer déjà que $\mathcal F'$ est $f'$-plat.

D'autre part, en tenant compte de (\textbf I, 3.6.5), (\textbf{II}, 5.4.2) et
(2.1.4), on peut remplacer $Y$ par $\operatorname{Spec}(\mathcal O_y)$,
autrement dit supposer que $Y=\operatorname{Spec}(A)$, où $A$ est un anneau
local. Sous ces conditions, on va prouver que $\mathcal F$ est $f$-plat. En
vertu de (5.13.3), $A$ est limite inductive de sous-anneaux locaux noethériens
$A_\lambda$ tels que l'injection canonique $A_\lambda\to A$ soit un
homomorphisme local. En vertu de (8.9.1), on peut supposer que
$X=X_\lambda\otimes_{A_\lambda}A$, $f=f_\lambda\otimes1_A$,
$\mathcal F=\mathcal F_\lambda\otimes_{A_\lambda}A$, pour un $\lambda$
convenable, $f_\lambda:X_\lambda\to Y_\lambda=\operatorname{Spec}(A_\lambda)$
étant un morphisme de type fini et $\mathcal F_\lambda$ un
$\mathcal O_{X_\lambda}$-Module cohérent. De même, on peut supposer que
$Y'=Y'_\lambda\otimes_{A_\lambda}A$, $g=g_\lambda\otimes1_A$, où
$g_\lambda:Y'_\lambda\to Y_\lambda$ est un morphisme de présentation finie~;
en outre, (8.10.5, (xii)) permet de supposer $g_\lambda$ propre. Comme par
hypothèse l'homomorphisme $A\to\Gamma(Y',\mathcal O_{Y'})$ est injectif et que
$A_\lambda$ est un sous-anneau de $A$, l'homomorphisme
$A_\lambda\to\Gamma(Y',\mathcal O_{Y'})$ est aussi injectif. Enfin, en vertu de
(11.2.6), on peut supposer $\lambda$ pris tel que
$\mathcal F'_\lambda=\mathcal F_\lambda\otimes_{Y_\lambda}Y'_\lambda$ soit
$f'_\lambda$-plat, puisque $\mathcal F'=\mathcal F'_\lambda\otimes_{Y'_\lambda}Y'$.
Ces remarques prouvent que l'on peut désormais supposer l'anneau $A$
noethérien, les autres hypothèses de (11.5.3) étant vérifiées.

Posons alors $B=\widehat A$, $Z=\operatorname{Spec}(B)$~; comme $B$ est un
$A$-module fidèlement plat $(\mathbf 0_{\mathrm I},7.3.5)$, il revient au même
de dire que $\mathcal F$ est $f$-plat ou que $\mathcal F\otimes_Y Z$ est
$Z$-plat (2.1.4). De même, si l'on pose $Z'=Y'\times_Y Z$, le morphisme
$Z'\to Y'$ est fidèlement plat (2.2.13), donc il revient au même de dire que
$\mathcal F'$ est $f'$-plat ou que $\mathcal F'\otimes_{Y'}Z'$ est $Z'$-plat~;
enfin $h=g_{(Z)}:Z'\to Z$ est propre (\textbf{II}, 5.4.2) et de type fini
(1.5.2), $\widehat A$ est noethérien, et si $z$ est son point
